\scanpage{26}

\uline{Lemme} $1.2.\deleted{\unsure{2|3|1|7}.}9$. \uline{Soient} $S$ \uline{un schéma local artinien}, \inserted{$\bar s$ un point géométrique de $S$}, $f:X\longrightarrow S$
\uline{un morphisme propre}\inserted{\uline{et lisse}}\deleted{\unsure{et plat}, de fibres \unsure{géométriques} \deleted{\illegible{1}} \unsure{réduites|connexes}}, \uline{et}
$G$ \uline{un schéma en groupes commutatif lisse de type fini sur} $S$.
\uline{Si} $H^0(X_{\bar s},\unsure{O|\mathcal O})=k(\bar s)$ \uline{et} $\underline{\mathrm{Pic}}^{\tau}X_{\bar s}=0$, \uline{et}
\deleted{alors le foncteur} \uline{si} $G_{\bar s}$ \uline{ne contient pas de sous-groupe} ${\mathbb G}_m$, \uline{alors le} \deleted{le} \uline{foncteur} $f^*$
\uline{de la catégorie des} $G$-\uline{torseurs sur} $S$ \uline{dans celle des} $f^*G$-\uline{torseurs sur} $X$
\uline{est une équivalence.}

% a. below the dividing line; its arrow ends immediately above b.
\inserted{a. Le problème posé est local sur $S$ pour la topologie étale ;
on se ramène donc au cas $S$ strictement hensélien, donc au cas où
\deleted{on peut donc supposer que} $X/S$ admet une section ; le foncteur $f^*$
est alors pleinement fidèle d'après \deleted{$1.2.2.4$}. $1.2.3$ et $(1.2.4.2)$}

b. \uline{Si} $S$ \uline{est le spectre d'un corps}\inserted{\deleted{séparablement clos}} et si $A$ \inserted{est} une variété
abélienne, alors $H^1(X,A)$ est de torsion (Raynaud [1, $\underline{\overline{\mathrm{XIII}}}$ $2.6$]),
donc réunion des images des groupes $H^1(X,{}_nA)$. \deleted{; ceux-ci
sont nuls par $1.2.2.7$}. Si $G$ est extension d'une variété
abélienne $A$ par un groupe \deleted{\unsure{$G_0$}} affine $G_0$, \deleted{$H^1(X,G)$ est donc}
% archive sheet number: 25

